
\section{The price must anticipate who chooses after conversion}
\label{sim:two-date-control-price}
The refinancing calendar tells us when an offer can settle. A second issue
arises when the security changes the person who chooses later company
policy. A price based on the old owner's future choice can be inconsistent
with a conversion that transfers control. The following finite example
isolates that dependency.

There are three dates: financing on 1 January 2027, a signal and policy
choice on 1 January 2028, and terminal distributions on 1 January 2029.
Investment costs 60 at the first date. Existing senior face is 40 and new
convertible face is 80. The gross discount factor is $21/20$ in each of the
two intervals. All cash recoveries occur at the terminal date. The equal
signal probabilities and the conditional probabilities below are a supplied
hypothetical pricing law, also used to evaluate expected cash utility.

Under the debt signal, the new claim remains debt and the old owner
chooses policy. Safe policy produces terminal estate 100. Risky policy
produces 20 or 220 with equal probabilities. Senior debt is paid first,
then the new claim up to its face, and old equity receives the residual.
The old owner consequently prefers risky expected cash 50 to safe cash
zero. The new creditor's expected cash under that choice is 40.

Under the conversion signal, the new claim is extinguished and its holder
receives 120 shares alongside 40 old shares, as in
Section~\ref{sim:conversion-book-rights}. The new holder now controls
policy. Safe estate is again 100; risky estate is 20 or 140 with equal
probabilities. After senior priority, safe policy gives the new holder
$3(100-40)/4=45$. Risky policy gives $75/2$, so the controlling holder
chooses safe.

To make the control distinction consequential, assume the old owner
receives a noncash private benefit worth 4 terminal utility units when
risky policy is chosen in this conversion branch. Old-owner risky cash
has expectation $25/2$, giving utility $25/2+4=33/2$, compared with safe
utility 15. Thus a stale old-controller assumption would select risky
policy. The benefit is deliberately hypothetical and never enters estate,
creditor recovery or accounting equity. At zero private benefit the old
owner also prefers safe policy, which checks the mechanism's limiting case.

The consistent continuation is therefore risky under the debt signal and
safe after conversion. Let $P$ denote date-zero proceeds paid by the new
holder. In this selected competitive-quote comparison, proceeds equal the
discounted expected cash the holder will receive:
\begin{equation}
 P=\frac{\tfrac12(40)+\tfrac12(45)}{(21/20)^2}
   =\frac{17000}{441},\qquad
 60-P=\frac{9460}{441}.
 \label{sim:control-dependent-quote}
\end{equation}
The two discount intervals matter: the signal date is not a payment date.
The second term is the cash the old owner must contribute at financing.
Available cash 20 is insufficient, while cash 22 is sufficient.

Affordability and willingness to invest are separate conditions. Under
this continuation, the old owner's expected terminal cash is
$(50+15)/2$. Its date-zero value is $13000/441$. Subtracting the
contribution gives $3540/441>0$, so participation exceeds the assumed
outside utility zero. That positive value does not supply the missing
liquid funds in the cash-20 case.

If old control were incorrectly retained after conversion, the new
holder's conversion-branch cash would instead be $75/2$. The corresponding
quote would be $15500/441$, requiring owner cash $10960/441$, more than
22. Charging the correct-control quote while executing that stale-control
policy leaves the lender short by $1500/441$. The discrepancy is a
failure to keep rights, continuation and price consistent. It does not
establish that conversion is economically superior.

Backward induction agrees with a separate enumeration of all four pure
strategies mapping the observed signal to a policy, including every local
deviation. A tie variant retains both optimal policies. Equal votes without
a supplied governance rule remain unresolved. Policies cannot depend on
the terminal realization, which is not known when the choice is made.

The owner's financing contribution is sunk by the policy date, so the
continuation ranking in this example is independent of the posted price.
This is a finite consistency calculation, not a general fixed-point or
market-clearing result. The typed conversion example starts at a separately
supplied fiscal opening; joining issuance, intervening measurement and final
liquidation into one accounting history remains a further composition task.
